\documentclass[11pt]{article}

\usepackage[margin=1in]{geometry}
\usepackage[T1]{fontenc}
\usepackage{lmodern}
\usepackage{microtype}
\usepackage{amsmath,amssymb}
\usepackage{booktabs,tabularx,array,longtable}
\usepackage{enumitem}
\usepackage{xcolor}
\usepackage{hyperref}
\usepackage{titlesec}
\usepackage{caption}
\usepackage{setspace}

\definecolor{headingblue}{RGB}{56,113,153}
\hypersetup{colorlinks=true,linkcolor=headingblue,citecolor=headingblue,urlcolor=headingblue}
\titleformat{\section}{\large\bfseries\color{headingblue}}{\thesection}{0.6em}{}
\titleformat{\subsection}{\normalsize\bfseries\color{headingblue}}{\thesubsection}{0.6em}{}
\titleformat{\subsubsection}{\normalsize\bfseries}{\thesubsubsection}{0.6em}{}
\captionsetup{font=small,labelfont=bf}
\setlength{\parindent}{1.2em}
\setlength{\parskip}{0pt}
\setlist[itemize]{leftmargin=1.5em,itemsep=0.25em,topsep=0.25em}
\setlist[enumerate]{leftmargin=1.8em,itemsep=0.45em,topsep=0.35em}
\renewcommand{\arraystretch}{1.05}
\newcommand{\ind}{\mathbb{I}}

\title{\color{headingblue}\bfseries Corridor Performance Analytics\\[0.25em]
Using Annotated LiDAR Data}
\author{Phase 1: Literature Review \& KPI Specification\\[0.5em]
\small Months 1--2 $\mid$ York University / University of Toronto + Ouster}
\date{June 2026}

\begin{document}
\maketitle

\begin{abstract}
This paper classifies the literature on corridor and intersection performance measurement and formalises Tier-1 Key Performance Indicators (KPIs) as functions of annotated LiDAR trajectories. Section 1 reviews manual, theoretical, detector-based, trajectory, and LiDAR sources across six KPI domains and identifies priority definitions for LiDAR transfer. Section 2 presents notation, governing equations with explicit variable definitions, and validation criteria aligned with Highway Capacity Manual and ATSPM practice. The specification provides the analytical foundation for real-time and daily corridor monitoring in subsequent project phases.
\end{abstract}

\section{Literature Review and Classification}

Signalized corridors remain the primary bottleneck in urban networks, yet agencies still rely heavily on periodic turning-movement counts, loop detectors, and manual studies to diagnose performance~\cite{hcm,bonneson}. Recent work with connected-vehicle trajectories and LiDAR-based tracking shows that individual vehicle paths can support continuous estimation of delay, queue extent, and signal utilization without repeated field visits~\cite{hunter,kotb,bhaskar}. A gap remains, however: most KPI definitions in manuals and agency practice assume aggregated detector data, not annotated object tracks with intermittent occlusion and identity switches.

The literature indicates that the most transferable KPI domains for a single
LiDAR-equipped intersection are delay and travel time, queue and storage,
signal utilization, and throughput. These measures have established nominal
definitions in the \emph{Highway Capacity Manual}, signal-timing manuals, and
Automated Traffic Signal Performance Measures (ATSPM) guidance, while
trajectory-based studies provide object-level computation procedures
\cite{nasem2022hcm,nasem2015stm,fhwa2020atspm,liu2009queue,
ashqer2022uav}. In contrast, corridor progression and multimodal or
transit-signal-priority measures often require multi-intersection trajectories,
signal-phase metadata, reliable modal classifications, or transit-event logs
that may not be available within a single LiDAR footprint
\cite{mahmud2022coordination,fhwaTSP}.

Phase 1 addresses this gap through a structured literature review (Tasks A1--A4) whose purpose is to (i) identify authoritative KPI definitions, (ii) classify evidence by source type and measurement domain, and (iii) determine which definitions transfer reliably to LiDAR-derived trajectories. The output is a traceability map from literature to the computable specification developed in Section~2. Evidence is organised into eight source classes (listed below); the domain taxonomy and priority matrix identify which definitions transfer to LiDAR trajectories.

\paragraph{Manuals (HCM, MUTCD, ITE):} Authoritative KPI definitions and units. Specifies control delay, queue length, saturation flow, and level of service. Project default when no LiDAR-specific rule applies.

\paragraph{Theory (Webster, Ak\c{c}elik, coordination):} Analytical models relating cycle length and arrivals to expected delay and progression. Provides validation ranges for LiDAR KPI magnitudes; does not define field measurement from trajectories.

\paragraph{Detectors (NCHRP, turning counts):} Documents how loops, radar, video, and manual counts validate delay, AOG, and volumes. External benchmarking targets for LiDAR-derived estimates.

\paragraph{Trajectories (CV, GPS, probe studies):} Per-vehicle KPI computation using rolling windows, integrals, and stop detection. Closest methodological precedent for LiDAR track aggregation and windowing.

\paragraph{LiDAR/vision (intersection point-cloud tracking):} Documents occlusion, ID switches, classification error, and track gaps. Defines LiDAR-specific adaptations when applying nominal HCM/ATSPM KPI definitions.

\paragraph{Multimodal/TSP (pedestrian, cyclist, transit priority):} Pedestrian delay, cyclist exposure, and TSP benefit metrics. Tier-2 scope in Phase 1: reviewed and scoped but not fully specified for LiDAR deployment.

\paragraph{Agency practice (ATSPM, FHWA):} Operational reporting thresholds, alert rules, and decision levers used by traffic engineers in daily operations.

\paragraph{Validation (trajectory QA):} Data-quality and benchmarking protocols. Consistency checks, coverage requirements, and spot-check methods to verify LiDAR-computed KPIs before publication to operators.

\subsection{Traceability and Definition-Transfer Rules}

The literature sources do not play identical roles in the KPI specification.
Manuals define the nominal meaning, units, and operational interpretation of a
KPI. Trajectory studies provide object-level computation and aggregation
logic. LiDAR and infrastructure-perception studies identify implementation
risks associated with occlusion, classification uncertainty, track
fragmentation, and identity switches. Detector, controller, and agency sources
provide external benchmarks and operational reporting conventions
\cite{nasem2022hcm,nasem2015stm,fhwa2020atspm,fhwa2006detectors,
kloeker2020multilidar,zimmer2023infradet3d}.

Each KPI is assigned to one of four transfer classes:

\begin{itemize}
    \item \textbf{Direct transfer:} the nominal definition is retained and all
    required variables are observable from trajectories and available signal
    metadata.

    \item \textbf{Adapted transfer:} the nominal definition is retained, but
    its implementation requires explicit treatment of occlusion, incomplete
    tracks, track merging, or limited sensor visibility.

    \item \textbf{Benchmark only:} the source provides expected ranges,
    external comparisons, or causal interpretation rather than a direct
    LiDAR-computation procedure.

    \item \textbf{Tier-2:} the KPI remains within the project scope, but its
    complete operational specification is deferred until the required
    corridor-wide, modal, or signal-event metadata are available.
\end{itemize}

The resulting traceability rule is therefore to use authoritative manuals for
KPI names, definitions, and units; trajectory studies for per-object
calculation logic; LiDAR and vision studies for data-quality adaptations; and
detector or agency guidance for validation and operational reporting.

\subsection{LiDAR Data Readiness and Scene Geometry}

Before computing any KPI, annotated LiDAR trajectories must be converted into
an operations-ready data structure. At minimum, every trajectory record should
contain a timestamp, object identifier, object class, position in world
coordinates, speed, heading, lane-group assignment, and a tracking-confidence
or continuity indicator. The static scene representation should contain lane
polygons, lane centerlines, stop bars, count lines, crosswalks, queue regions,
turn-bay storage lengths, waiting areas, and the boundary of the monitored
intersection influence area
\cite{fhwa2006detectors,zimmer2023infradet3d,nayak2025dual}.

LiDAR observations should not be treated as error-free ground truth by
default. Published KPIs may be biased by timestamp misalignment, spatial
calibration errors, blind zones, classification errors, identity switches, and
fragmented trajectories. Consequently, every reported KPI should be
accompanied by quality metadata, including the number of valid tracks, the
percentage of the reporting window with acceptable coverage, the number of
track merges, and a visibility or confidence flag
\cite{kloeker2020multilidar,shang2026roadside}.

Corridor KPIs are grouped into six domains (Table~\ref{tab:domains}). Domains D1--D4 (delay, queue, signal utilization, throughput) are well established in HCM and ATSPM practice and are feasible from LiDAR tracks within a single intersection footprint. Domains D5--D6 (progression, multimodal) are retained in the review but scoped as Tier-2 in Phase 1 because they often require corridor-wide or mode-specific metadata beyond a single sensor node.

\begin{table}[htbp]
\centering
\caption{KPI measurement domains}
\label{tab:domains}
\small
\begin{tabularx}{\textwidth}{@{}p{0.10\textwidth}X p{0.16\textwidth}p{0.17\textwidth}@{}}
\toprule
Domain & Representative KPIs & Cadence & LiDAR feasibility \\
\midrule
D1 Delay \& travel time & Control delay, stopped delay, segment travel time & RT / Daily & High \\
D2 Queue \& storage & Queue length, spillback, storage utilization & RT / Daily & High \\
D3 Signal utilization & AOG, split failure, platoon ratio & RT / Daily & Medium--High \\
D4 Throughput & Approach volume, discharge rate, turning volumes & RT / Daily & High \\
D5 Progression & Stop frequency, offset effectiveness & Daily & Medium \\
D6 Multimodal & Pedestrian delay, TSP benefit & Daily / Event & Medium--Low \\
\bottomrule
\end{tabularx}
\end{table}

Table~\ref{tab:priority} cross-classifies source classes against domains. The matrix shows that Classes A, D, and G dominate D1--D4, while Class E is essential for documenting LiDAR-specific adaptations (track gaps, classification error) even when the nominal KPI definition follows the HCM.

\begin{table}[htbp]
\centering
\caption{Source class $\times$ KPI domain priority matrix. P (priority): this source class receives full definitional treatment in the KPI specification for that domain. S (secondary): the source is referenced with scope notes only, not a complete specification. n/a: not applicable in Phase 1.}
\label{tab:priority}
\scriptsize
\begin{tabular}{@{}lcccccc@{}}
\toprule
Source class & Delay & Queue & Signal util. & Throughput & Progression & Multimodal \\
\midrule
A. Manuals & P & P & P & P & P & S \\
B. Theory & P & P & P & P & P & n/a \\
C. Detectors & S & S & P & P & S & S \\
D. Trajectories & P & P & P & P & P & S \\
E. LiDAR/vision & P & P & S & P & S & S \\
F. Multimodal/TSP & S & S & n/a & S & n/a & P \\
G. Agency & P & P & P & P & P & S \\
H. Validation & P & P & P & P & P & P \\
\bottomrule
\end{tabular}
\end{table}

\section{KPI Specification for LiDAR Trajectories}

Building on the classification in Section~1, this section formalises Tier-1 corridor KPIs as LiDAR-computable functions of annotated trajectories. Definitions follow HCM and ATSPM conventions~\cite{hcm,smaglik} where applicable; trajectory-specific terms (stop flag, approach polygon, track continuity) follow probe and LiDAR precedents~\cite{bhaskar,kotb}. Table~\ref{tab:tier1} summarises the seven Tier-1 KPIs; definitions are given below. Subsequent subsections introduce notation and present the governing equations.

\begin{table}[htbp]
\centering
\caption{Tier-1 KPIs and reporting cadence}
\label{tab:tier1}
\small
\begin{tabular}{@{}llll@{}}
\toprule
KPI & Domain & Units & RT / Daily \\
\midrule
Control delay $\bar d_a$ & D1 & s/veh & Both \\
Stopped delay $d_{i,\mathrm{stop}}$ & D1 & s/veh & Both \\
Queue length $Q_{a,95}$ & D2 & m & Both \\
Arrivals on green (AOG) & D3 & \% & Both \\
Split failure rate (SF) & D3 & \% of cycles & Both \\
Approach volume $V_a$ & D4 & veh/h & Both \\
Discharge rate $q_a$ & D4 & veh/h/lane & RT \\
\bottomrule
\end{tabular}
\end{table}

\subsection{Tier-1 KPI Definitions}

\begin{enumerate}
\item \textbf{Control delay ($\bar d_a$; s/veh):} Mean excess travel time per vehicle on an approach relative to free-flow conditions, from approach entry to the stop line. Reported as the arithmetic mean over valid traversals in the analysis window (RT rolling window or daily calendar window).

\item \textbf{Stopped delay ($d_{i,\mathrm{stop}}$; s/veh):} Time a vehicle spends at crawl or standstill (speed $\leq v_{\mathrm{th}}$, typically 2--3 m/s) while traversing the approach. Computed by integrating the stop indicator over the vehicle's traversal interval.

\item \textbf{Queue length, 95th percentile ($Q_{a,95}$; m):} Farthest upstream distance among stopped vehicles from the stop line at each time step; the operational KPI is the 95th percentile over the reporting window to reduce sensitivity to single-frame outliers. Set to zero when no stopped vehicles are present.

\item \textbf{Arrivals on green (AOG; \%):} Share of vehicles arriving at the stop line during effective green. A vehicle counts as an arrival when within $\varepsilon \approx 3$ m of the stop line; green membership requires the controller phase at arrival time to match effective green $g_{\mathrm{eff}}$.

\item \textbf{Split failure rate (SF; \% of cycles):} Share of signal cycles in which queue remains at the stop line when green ends ($Q_a(t_{c,\mathrm{end}})>0$). Indicates insufficient green time relative to demand.

\item \textbf{Approach volume ($V_a$; veh/h):} Number of vehicles entering the approach per hour, counting unique tracks with entry time $t_{i,\mathrm{entry}}$ in the analysis window.

\item \textbf{Discharge rate ($q_a$; veh/h/lane):} Vehicles departing through the stop line per hour of effective green per lane (stopped-to-moving transition during green). Values approaching HCM saturation flow indicate near-capacity operation. RT reporting only.
\end{enumerate}

\subsection{Notation and Data Model}

Let each vehicle track be indexed by $i\in\mathcal{T}$, where $\mathcal{T}$ is the set of valid trajectories in the analysis period. Track $i$ carries timestamped state $(t,\mathbf{x}_i(t),\mathbf{v}_i(t),c_i)$, where $\mathbf{x}_i(t)$ is position, $\mathbf{v}_i(t)$ velocity, and $c_i$ the object class (vehicle, pedestrian, etc.). Approach $a\in\mathcal{A}$ is defined by entry and stop-line polygons; only tracks with $c_i=\text{vehicle}$ and complete traversal of approach $a$ enter the vehicular KPI set.

\begin{table}[htbp]
\centering
\caption{Primary notation for KPI computation}
\label{tab:notation}
\small
\begin{tabularx}{\textwidth}{@{}p{0.16\textwidth}p{0.10\textwidth}X@{}}
\toprule
Symbol & Units & Definition \\
\midrule
$v_i(t)$ & m/s & Speed magnitude, $v_i(t)=\lVert\mathbf{v}_i(t)\rVert$ \\
$v_{\mathrm{th}}$ & m/s & Stopped-speed threshold (typically 2--3 m/s) \\
$S_i(t)$ & -- & Stop indicator, $S_i(t)=1$ if $v_i(t)\leq v_{\mathrm{th}}$ \\
$s_i(t)$ & m & Signed distance from track $i$ to stop line (negative upstream) \\
$t_{i,\mathrm{entry}}$ & s & Time track $i$ enters approach polygon of $a$ \\
$t_{i,\mathrm{stop}}$ & s & Time track $i$ first reaches stop line on approach $a$ \\
$t_{i,\mathrm{exit}}$ & s & Time track $i$ leaves approach or crossing zone \\
$L_a$ & m & Length of approach $a$ from entry to stop line \\
$v_{\mathrm{ff}}$ & m/s & Free-flow speed on approach $a$ (posted or calibrated) \\
$t_{i,\mathrm{ff}}$ & s & Free-flow traversal time, $L_a/v_{\mathrm{ff}}$ \\
$W_t$ & s & Rolling analysis window ending at time $t$ (length 60--300 s) \\
$W_{\mathrm{day}}$ & s & Calendar-day aggregation window \\
$N_a$ & -- & Number of valid vehicle traversals on approach $a$ in window \\
$g_{\mathrm{eff}}$ & s & Effective green for movement on approach $a$ \\
$t_{c,\mathrm{end}}$ & s & End of green for cycle $c$ \\
$N_{\mathrm{cycles}}$ & -- & Number of signal cycles in analysis window \\
$\tau$ & h & Duration of volume-count window \\
$T_g$ & h & Duration of green interval used for discharge \\
\bottomrule
\end{tabularx}
\end{table}

\subsection{Derived Features and Temporal Aggregation}

From raw LiDAR tracks we compute kinematic features used across KPIs:
\begin{equation}
 v_i(t)=\lVert\mathbf{v}_i(t)\rVert,\qquad
 S_i(t)=\ind\!\left[v_i(t)\leq v_{\mathrm{th}}\right],
\end{equation}
where $\ind[\cdot]$ is the indicator function. A track is excluded if any gap in observations exceeds $\Delta t_{\max}$ (project default: 3 s) within $[t_{i,\mathrm{entry}},t_{i,\mathrm{exit}}]$.

KPI values are aggregated over rolling windows for real-time (RT) monitoring and fixed windows for daily summaries:
\begin{equation}
 \bar K_{\mathrm{RT}}(t)=\frac{1}{\lvert\mathcal{I}_{W_t}\rvert}
 \sum_{i\in\mathcal{I}_{W_t}}k_i,
 \qquad
 \bar K_{\mathrm{daily}}=\frac{1}{\lvert\mathcal{I}_{W_{\mathrm{day}}}\rvert}
 \sum_{i\in\mathcal{I}_{W_{\mathrm{day}}}}k_i,
\end{equation}
where $k_i$ is the per-vehicle KPI contribution (e.g., $d_i$ for delay), $\mathcal{I}_{W_t}$ is the set of valid trajectories with $t_{i,\mathrm{stop}}\in W_t$, and $\lvert\cdot\rvert$ denotes set cardinality. RT windows slide continuously; daily windows use non-overlapping local calendar days.

A fragmented trajectory may be merged with a subsequent fragment only when
the fragments have consistent object classes and lane assignments, a short
temporal separation, and feasible kinematic continuity between the end of the
first fragment and the beginning of the second. If multiple candidate matches
exist, the merge should be rejected unless one assignment is clearly superior.
Every accepted merge should be retained in an audit log
\cite{kloeker2020multilidar,shang2026roadside}.

Tracks beginning or ending inside the monitored analysis region are considered
censored. They may contribute to KPIs that require only an observed event,
such as a count-line crossing, but should not contribute to full-approach
travel time or control delay unless both required boundaries are observed.

\subsection{Delay Metrics (Domain D1)}

\paragraph{Control delay.} Following HCM control-delay concepts~\cite{hcm}, the excess time for vehicle $i$ on approach $a$ relative to free-flow conditions is
\begin{equation}
 d_i=
 \underbrace{\left(t_{i,\mathrm{stop}}-t_{i,\mathrm{entry}}\right)}_{\text{actual traversal time}}
 -\underbrace{t_{i,\mathrm{ff}}}_{\text{free-flow time}},
 \qquad
 t_{i,\mathrm{ff}}=\frac{L_a}{v_{\mathrm{ff}}}.
\end{equation}
The approach-level control delay reported to operators is the arithmetic mean over valid traversals:
\begin{equation}
 \bar d_a=\frac{1}{N_a}\sum_{i=1}^{N_a}d_i,
 \qquad
 N_a=\lvert\mathcal{I}_a\cap\mathcal{I}_W\rvert.
\end{equation}

Equation (4) should be applied only to trajectories for which both approach
entry and stop-line arrival are observed. If the LiDAR trajectory begins or
ends within the approach polygon, the calculated value should be labelled
\emph{observed-segment delay} rather than full-approach control delay. Such
censored observations should be reported separately and excluded from the
primary approach mean unless a validated correction procedure is applied
\cite{wang2024state}.

\paragraph{Stopped delay.} Stopped delay isolates time spent at crawl or standstill, computed by integrating the stop indicator:
\begin{equation}
 d_{i,\mathrm{stop}}=\int_{t_{i,\mathrm{entry}}}^{t_{i,\mathrm{exit}}}S_i(t)\,dt.
\end{equation}
In discrete LiDAR samples at interval $\delta t$, (5) is evaluated as $d_{i,\mathrm{stop}}\approx\sum_n S_i(t_n)\,\delta t$.

\subsection{Queue Length (Domain D2)}

Queue extent at time $t$ is the maximum upstream distance among stopped vehicles on approach $a$:
\begin{equation}
 Q_a(t)=\max\left\{\lvert s_i(t)\rvert:S_i(t)=1,\ i\in\mathcal{T}_a\right\},
\end{equation}
where $\mathcal{T}_a$ is the set of tracks currently on approach $a$. The operational KPI is the 95th percentile over reporting window $W$ to reduce sensitivity to single-frame outliers:
\begin{equation}
 Q_{a,95}=P_{95}\left\{Q_a(t):t\in W\right\}.
\end{equation}
If no stopped vehicles are present, $Q_a(t)=0$.

The stopped-speed condition alone is insufficient to establish queue
membership because an isolated turning, parking, or yielding vehicle may also
move below $v_{\mathrm{th}}$. A vehicle should therefore be classified as a
queue member only when it is lane-consistent, located within the designated
queue region, and has plausible spacing from another queued vehicle or from
the stop line. The queue estimate should also be accompanied by a visibility
flag because occlusion of the upstream queue tail can systematically
underestimate queue extent
\cite{liu2009queue,ashqer2022uav,kloeker2020multilidar}.

\subsection{Signal Utilization (Domain D3)}

\paragraph{Arrivals on green (AOG).} Let $N_{\mathrm{arr}}$ be vehicles whose stop-line arrival time $t_{i,\mathrm{stop}}$ falls in window $W$, and $N_{\mathrm{green}}$ the subset arriving during effective green $g_{\mathrm{eff}}$:
\begin{equation}
 \mathrm{AOG}=100\cdot\frac{N_{\mathrm{green}}}{N_{\mathrm{arr}}}.
\end{equation}
Vehicle $i$ is counted in $N_{\mathrm{arr}}$ when $\lvert s_i(t_{i,\mathrm{stop}})\rvert\leq\varepsilon$ with tolerance $\varepsilon\approx3$ m. Membership in $N_{\mathrm{green}}$ requires the controller phase at $t_{i,\mathrm{stop}}$ to match $g_{\mathrm{eff}}$, obtained from signal logs or Ouster timing metadata when available.

AOG should be considered operationally valid only when the LiDAR and
controller clocks have been explicitly synchronized and timestamp drift has
been evaluated. If signal state is inferred rather than obtained from
controller or SPaT logs, the resulting AOG value should be labelled
provisional and reported with the inference method and associated uncertainty
\cite{fhwa2020atspm}.

\paragraph{Split failure.} For signal cycle $c$, define the split-failure indicator
\begin{equation}
 F_c=\ind\!\left[Q_a(t_{c,\mathrm{end}})>0\right],
\end{equation}
i.e., queue remains at the stop line when green ends. The split-failure rate over $N_{\mathrm{cycles}}$ observed cycles is
\begin{equation}
 \mathrm{SF}=100\cdot\frac{1}{N_{\mathrm{cycles}}}
 \sum_{c=1}^{N_{\mathrm{cycles}}}F_c.
\end{equation}

The residual-queue test should use the queue-membership definition introduced
in Domain D2 rather than the presence of any isolated stopped vehicle.
Cycle-level split-failure results should also report the number of complete
cycles and the percentage of cycles with adequate approach visibility.

\subsection{Throughput (Domain D4)}

\paragraph{Approach volume.} The hourly entry flow on approach $a$ over a window of duration $\tau$ (hours) is
\begin{equation}
 V_a=\frac{N_{a,\mathrm{entry}}}{\tau},
\end{equation}
where $N_{a,\mathrm{entry}}$ counts unique vehicle tracks with $t_{i,\mathrm{entry}}\in W$.

Because one physical vehicle may appear under multiple track identifiers,
approach volume must be based on unique physical trajectories after
continuity-based deduplication. Count-line crossings should also satisfy a
direction condition to prevent repeated counts when an object oscillates near
the line
\cite{nayak2025dual,shang2026roadside}.

\paragraph{Discharge rate.} During a green interval of duration $T_g$ (hours), the discharge rate is
\begin{equation}
 q_a=\frac{N_{a,\mathrm{depart}}}{T_g},
\end{equation}
where $N_{a,\mathrm{depart}}$ is the number of vehicles crossing the stop line while $S_i$ transitions from stopped to moving during green. Values of $q_a$ approaching saturation flow $s$ (veh/h/lane) from HCM~\cite{hcm} indicate near-capacity operation.

When a standing queue exists at the beginning of green, saturation discharge
may also be estimated from stable post-start-up headways. If
$h_4,\ldots,h_k$ denote the headways after the initial start-up vehicles, then

\begin{equation}
 s_a=\frac{3600}{\overline{h}_{4:k}},
\end{equation}

where $s_a$ is the observed saturation-flow estimate in veh/h/lane. This
headway-based estimate should be reported only for cycles containing a
standing queue at the beginning of effective green
\cite{nasem2022hcm,webster1958,akcelik1981signals}.

\subsection{Validation}

Published KPIs must satisfy three checks. \textbf{Internal consistency:} $N_{\mathrm{green}}+N_{\mathrm{red}}\approx N_{\mathrm{arr}}$; $\bar d_a$ positively correlated with $Q_{a,95}$; turning-movement volumes sum to approach entry volume within $\pm5\%$. \textbf{External benchmark:} where loop counts or controller AOG logs exist, compare $V_a$ and AOG (targets: $\pm10\%$ volume, $\pm5\%$ AOG). \textbf{Data quality:} require $N_a\geq15$ per RT window, lane coverage $\geq80\%$, and exclude tracks with gaps $>\Delta t_{\max}$.

The Phase 1 deliverable documents Equations (1)--(13) with the notation in Table~\ref{tab:notation}, forming the specification handed to Phase 2 KPI computation (Tasks T2.1--T2.2).



\appendix
\section*{Appendix: Phase 1 Execution Roadmap}
\addcontentsline{toc}{section}{Appendix: Phase 1 Execution Roadmap}

The following material supports project management and is secondary to the analytical content in Sections 1--2.

\subsection*{Stakeholder Interviews (Tasks C1--C4)}
\begin{itemize}
\item \textbf{C1 Identify cities:} Minimum three agencies, mixed size and region; target traffic operations and corridor management staff.
\item \textbf{C2 Interview questions:} KPI priorities, reporting cadence, thresholds used in practice, remedies commonly applied (signal timing, coordination, phasing), and data trust requirements.
\item \textbf{C3 Conduct interviews:} Semi-structured sessions; document current ATSPM or arterial analytics workflows.
\item \textbf{C4 Synthesize findings:} Cross-city comparison table; reconcile with literature classification; flag KPIs for Phase 2 priority implementation.
\end{itemize}

\subsection*{Schedule and Milestones}
\begin{table}[htbp]
\centering
\caption{Phase 1 work breakdown aligned to SOW tasks T1.1--T1.10}
\label{tab:workbreakdown}
\small
\begin{tabularx}{\textwidth}{@{}p{0.22\textwidth}p{0.16\textwidth}X@{}}
\toprule
Stream & Task IDs & Activities \\
\midrule
Data engineering & T1.1--T1.3 & LiDAR ingestion pipeline; QA/QC scripts; preliminary features (speed, acceleration, stops) \\
Literature \& KPI & T1.4--T1.6 & Classified literature review; KPI formulas and cadence; feasibility validation on sample data \\
Analytics \& diagnostics & T1.7--T1.10 & City interviews; exploratory trajectory analysis; KPI stability evaluation; preliminary anomaly signatures \\
\bottomrule
\end{tabularx}
\end{table}

\subsection*{Deliverables and Phase 2 Handoff}

Phase 2 implements rolling-window KPI computation (Tasks T2.1--T2.9) using the specification in Section~2. Milestones: Week 4 (literature memo draft, interviews underway); Week 8 (KPI specification, validation plan, stakeholder synthesis, Phase 2 handoff).

\begin{table}[htbp]
\centering
\caption{Phase 1 deliverables and downstream use}
\label{tab:deliverables}
\small
\begin{tabularx}{\textwidth}{@{}p{0.25\textwidth}X p{0.25\textwidth}@{}}
\toprule
Deliverable & Contents & Feeds \\
\midrule
Best-practices memo & Standard KPI definitions; source-class synthesis; domain priority matrix & KPI spec; agency alignment \\
KPI specification draft & LiDAR-ready formulas, units, cadence, data requirements per KPI & Phase 2 real-time engine \\
Validation \& benchmarking plan & Consistency checks, timing-plan comparison, spot-check protocol & Phase 2 validation (T2.4--T2.5) \\
Stakeholder interview synthesis & Findings from $\geq3$ cities; priority KPIs and decision levers & Phase 2 reporting format; Phase 3 diagnostics \\
Data pipeline \& QA baseline & Ingestion, profiling scripts, preliminary features & KPI feasibility (T1.6); Phase 2 automation \\
\bottomrule
\end{tabularx}
\end{table}

\subsection*{Acknowledgments}

This roadmap implements Phase 1 of the Statement of Work for the Corridor Performance Analytics project, a collaboration between York University / University of Toronto and Ouster.

\bibliographystyle{unsrt}
\bibliography{corridor_lidar_lit}

\end{document}
